\section{Full multiplicity from transverse jets along a binary plane}

\subsection{O-020: arbitrary second rows with third row at most two}
Let $r\geq3$ and $n\geq r-1$. For every outer degree $d$
and every partition $\lambda\vdash nd$ with
$\ell(\lambda)\leq r$ and $\lambda_3\leq2$,
\[
                         K_n(\lambda)=a_\lambda(d[n]).
                                                               \tag{O-020.1}
\]
Missing third parts are zero. There is no bound on
$\lambda_2$ and no upper bound on $d$.
The geometric input is dominance of the determinant projection
onto the entire coefficient space of degree-$n$ forms whose
total degree outside a fixed binary plane is at most two.
Higher transverse coefficients remain unrestricted.

Write the variables as $x,y,z_1,\ldots,z_R$, where $R=r-2$.
The target coefficient space is
\[
 \mathcal T_{n,2}=
 \operatorname{Sym}^n\mathbb C^2
 \oplus\left(\operatorname{Sym}^{n-1}\mathbb C^2
                               \otimes\mathbb C^R\right)
 \oplus\left(\operatorname{Sym}^{n-2}\mathbb C^2
                               \otimes\operatorname{Sym}^2\mathbb C^R\right).
                                                               \tag{O-020.2}
\]
Its dimension is $(n+1)+Rn+\binom{R+1}{2}(n-1)$.
The map from all $n$-by-$n$ linear pencils in these variables
to $\mathcal T_{n,2}$, obtained by deleting terms of
total $z$-degree at least three from the determinant,
has a surjective differential at an explicit integral pencil.

\paragraph{A difference-labeled upper-triangular pencil.}
Since $n\geq R+1$, put
\[
 \zeta_i=z_i\quad(1\leq i\leq R),\qquad
 \zeta_i=0\quad(R+1\leq i\leq n),\qquad
 u_{ij}=\zeta_i-\zeta_j\quad(i<j).
\]
Let $U=(u_{ij})$ be strictly upper triangular, choose distinct integers
$a_1,\ldots,a_n$, and define
\[
 \ell_i=y-a_ix,\qquad \Delta=\prod_{i=1}^n\ell_i,
             A_0=\operatorname{diag}(\ell_1,\ldots,\ell_n)-U.
                                                               \tag{O-020.3}
\]
Then $\det A_0=\Delta$.

\paragraph{Binary coefficients and residue coordinates.}
The $n$ polynomials $\Delta/\ell_i$ form a basis of the
binary forms of degree $n-1$. After setting $x=1$, their
evaluation matrix at $y=a_j$ is diagonal with nonzero
diagonal entries.
Send a polynomial $p(y)$ of degree
at most $n-2$ to the vector
\[
                \left(\frac{p(a_i)}{\Delta'(a_i)}\right)_{i=1}^n.
\]
Lagrange interpolation shows that its coordinates sum to
zero; the map is injective and both spaces have dimension
$n-1$, so it identifies this polynomial space with the
zero-sum hyperplane $H_n\subseteq\mathbb C^n$. The image of
$\Delta/((y-a_i)(y-a_j))$ is
$(e_i-e_j)/(a_i-a_j)$. Thus, after rescaling each independent
edge parameter by a nonzero scalar, the binary coefficients
$\Delta/(\ell_i\ell_j)$ can be treated as the vectors
$e_i-e_j$ in $H_n$.

\paragraph{A quadratic incidence isomorphism.}
First take $n=R+1$, and let $Z_1$ denote the linear forms
in $z_1,\ldots,z_R$. Consider
\[
 \Phi:\bigoplus_{1\leq i<j\leq R+1} Z_1
       \longrightarrow H_{R+1}\otimes\operatorname{Sym}^2Z_1,
 \qquad
 (f_{ij})\longmapsto
   \sum_{i<j}(e_i-e_j)(\zeta_i-\zeta_j)f_{ij}.
                                                               \tag{O-020.4}
\]
Its source and target both have dimension
$R\binom{R+1}{2}$. The map is an isomorphism.

To prove injectivity, write $f_{ji}=f_{ij}$ and
$f_{ij}=\sum_{k=1}^R c_{ij,k}z_k$. A kernel element satisfies
\[
          \sum_{j\neq i}(\zeta_i-\zeta_j)f_{ij}=0
                                  \quad(1\leq i\leq R+1).
\]
For distinct $i,j\leq R$, the coefficient of $z_j^2$
in the $i$-th equation gives $c_{ij,j}=0$.
The $j$-th equation likewise gives $c_{ij,i}=0$.
For three distinct $i,j,k\leq R$, the coefficient of
$z_jz_k$ in the $i$-th equation gives
$c_{ij,k}+c_{ik,j}=0$. Applying the same relation at
$j$ and $k$, and using $c_{ji,k}=c_{ij,k}$, shows that
$c_{ij,k}=-c_{ij,k}$. Hence it is zero in characteristic
zero. This also covers $R=1,2$, where there are no such
remaining triple coefficients. Every $f_{ij}$ with
$i,j\leq R$ therefore vanishes. The $i$-th equation now
reduces to $z_i f_{i,R+1}=0$, which forces the remaining
linear forms to vanish. This proves injectivity and the
dimension count proves surjectivity.

For $n>R+1$ add the vertices one at a time. The previous
image is $H_{n-1}\otimes\operatorname{Sym}^2Z_1$, embedded
with last coordinate zero. For a new vertex $v$, the
edges $\{i,v\}$, $1\leq i\leq R$, have labels $z_i$.
Their new residue coordinates are $-z_i f_{iv}$, which
span every quadratic polynomial, since every quadratic
monomial is a product of a variable with a linear form.
Choose these parameters to obtain any prescribed new
coordinate. The remaining error has new coordinate zero
and total coordinate sum zero, so it belongs to the old
image and can be removed by the old edge parameters.
This proves surjectivity onto
$H_n\otimes\operatorname{Sym}^2Z_1$ for every $n\geq R+1$.
It checks every additional residue coordinate, rather than
merely comparing the numbers of parameters.

\begin{proof}[Proof of the surjective differential]
The cofactor derivative formula is
\[
             d(\det)_{A_0}(H E_{ij})
                         =H\,\Delta\,(A_0^{-1})_{ji}.
                                                               \tag{O-020.5}
\]
The expression is polynomial although it is convenient to
compute the inverse over $\mathbb C(x,y,z)$.
Since $U$ is strictly upper triangular, an inverse entry
is a finite sum over increasing paths. In particular,
\[
 (A_0^{-1})_{ii}=\ell_i^{-1},\qquad
 (A_0^{-1})_{ij}
       =\frac{u_{ij}}{\ell_i\ell_j}
                    +\text{terms of total $z$-degree at least two}
                       \quad(i<j).                    \tag{O-020.6}
\]
Every extra path edge contributes another $z$ factor.
Multiplication by $\Delta$ clears the denominators,
because a strictly increasing path visits each vertex once.

Vary diagonal entries first. Taking $H=x$ or $H=y$
in (O-020.5) spans every binary degree-$n$ form:
the polynomials $\Delta/\ell_i$ span degree $n-1$,
and multiplication by $x,y$ spans the next degree.
Taking $H=z_a$ on diagonal entries spans independently
every term linear in the $z$ variables with an arbitrary
binary coefficient of degree $n-1$.

For any edge $\{i,j\}$, $i<j$, vary the opposite,
lower-triangular entry by an arbitrary $H\in Z_1$.
Equations (O-020.5)--(O-020.6), after deleting transverse
degree at least three, give exactly
\[
                  (\zeta_i-\zeta_j)H
                          \frac{\Delta}{\ell_i\ell_j}.
                                                               \tag{O-020.7}
\]
The terms from longer paths have degree at least three
after multiplication by $H$ and disappear from the
target projection. The residue identification and the
surjectivity of (O-020.4), including its added-vertex
extension, show that these directions span every quadratic
transverse polynomial with arbitrary binary coefficients
of degree $n-2$.

These directions therefore span every summand of
(O-020.2). Diagonal directions control transverse orders
zero and one, while the lower-triangular directions just
used have no contributions in those orders. All matrix-entry coefficients are
independent parameters. Thus the differential is surjective,
with no requirement to fix higher transverse terms.
The algebraic lowest-homogeneous-term argument in M-013
now proves dominance of this coefficient projection.
\end{proof}

\begin{proof}[Proof of the multiplicity equality]
Let $w$ be any degree-$d$ highest-weight vector of weight
$\lambda$ in the determinant ideal. It depends only on
the first $r$ coordinates. By the $j=3$ support bound
of F-016, every coefficient generator used by $w$ has
total degree at most $\lambda_3\leq2$ in variables
$3,\ldots,r$. Hence $w$ is a polynomial in the
coordinates of $\mathcal T_{n,2}$.
Every linear pencil determinant belongs to $\Omega_n$,
so this polynomial vanishes after the dominant projection
just proved. Its pullback is injective, so $w=0$.
This holds for every highest-weight copy, and therefore
the ideal multiplicity space vanishes entirely. The
ambient quotient then gives (O-020.1).
\end{proof}

For permanent shapes $r\leq m^2$, the condition
$n\geq m^2-1$ is enough for this full equality whenever
$\lambda_3\leq2$, with the one- and two-row cases already
covered by O-006. This conclusion is stronger than a
comparison with a particular permanent multiplicity: every
ambient copy survives, even with an arbitrarily large second
row. It is a different sufficient condition from M-013's
bounded second-row jet theorem.

\paragraph{An independently checkable tree construction.}
The earlier proof with the weaker sufficient bound
$n\geq2R=2r-4$ is retained in
\texttt{binary\_transverse\_jets\_tree.tex}.
It colors $R$ edge-disjoint spanning trees by the individual
variables $z_a$; each tree's edge cofactors form a binary
degree-$n-2$ basis. The paths
$j,(j-i,j+i)_{1\leq i\leq R-1},j-R$ on $\mathbb Z/(2R)$
give an explicit packing, including $R=1$.
Its independent cofactor calculations pass eight exact
Jacobian rank tests, including the smallest case $R=1,n=2$.
The difference-label argument above improves the threshold
by using the quadratic incidence isomorphism instead of
disjoint color classes.

\paragraph{The exact integral incidence cokernel.}
The particular map $\Phi$ in (O-020.4), with $n=R+1$,
is defined between free abelian groups of the same rank.
Its integral cokernel is
\[
          \operatorname{coker}(\Phi_{\mathbb Z})
                      \cong(\mathbb Z/2\mathbb Z)^{\binom R3}.
                                                               \tag{O-020.8}
\]
Here the target uses the ordinary monomial basis of
$\operatorname{Sym}^2\mathbb Z^R$. Delete the last residue
coordinate to identify $H_{R+1,\mathbb Z}$ with $\mathbb Z^R$.
The $R^2$ coefficients of the edges from $i$ to the zero
vertex occur with coefficient one in the rows given by
vertex $i$ and a monomial $z_i z_k$. They occur in no
other rows after the last residue coordinate is deleted.
Unimodular elimination removes these columns and rows.

Next, the two endpoint coefficients of each finite edge
$\{i,j\}$ occur with coefficient minus one in the remaining
rows $(i,z_j^2)$ and $(j,z_i^2)$, respectively.
They are unit pivots and can also be eliminated.
The remaining coordinates split by unordered triples
$\{i,j,k\}$. Each triple has the three variables
$c_{ij,k},c_{ik,j},c_{jk,i}$ and the three relations
\[
             a+b,\qquad a+c,\qquad b+c,
\]
up to signs. This $3$-by-$3$ block has Smith form
$\operatorname{diag}(1,1,2)$: a unit entry and a unit
$2$-by-$2$ minor give its first two invariant factors,
and its determinant has absolute value two. This proves
(O-020.8), including $R=1,2$, where there are no triples.
Thus this selected quadratic incidence submap is an
isomorphism after inverting two, and its rank drops by
$\binom R3$ after reduction modulo two.
No failure of the full determinant jet map in characteristic
two is inferred; other coefficient directions are not
included in $\Phi$.

\paragraph{Geometric dictionary.}
The target is the degree-$n$ part of the quotient by
$(z_1,\ldots,z_R)^3$, the second infinitesimal neighborhood
transverse to the binary coordinate plane. It is a finite
locally free coefficient module. The difference-labeled graph, diagonal
linear forms, and cofactor derivatives define integral
morphisms. Distinctness of the $a_i$ and the nonzero
Vandermonde factors and the precisely computed two-torsion
in (O-020.8) explain the passage to
$\mathbb Q$ or $\mathbb C$; no equal-rank assertion in
every residue characteristic is inferred.
The suffix support is a torus/comodule calculation under
SGA~3 I, Propositions~4.7.2--4.7.3, pp.~34--36,
and the Cauchy and Schur functors of Macdonald I(4.3),
p.~63, and Appendix~A(5.3)--(5.6), pp.~155--156.
Dominance injects the actual coordinate algebra of the
detected coefficients, which makes the whole highest-weight
ideal component vanish. Neither exact lifting of each
truncated form nor containment of that entire linear space
in the determinant closure is asserted.
